\section{Fundamental convex geometry}

  Let $V$ be a real vector space.

\subsection{Convex sets}

  \begin{definition}
    A set $K \subseteq V$ is called \emph{convex} if for any $x, y \in K$ and any $t \in [0, 1]$, we have
    \[t x + (1 - t) y \in K.\]
  \end{definition}

  \begin{example}
    \Yang{}
  \end{example}

  \begin{proposition}
    The intersection of any collection of convex sets is convex.
  \end{proposition}

  \begin{definition}
    Let $S$ be a subset of $V$.
    The \emph{convex hull} of $S$, denoted by $conv(S)$, is the smallest convex set containing $S$.
  \end{definition}

  \begin{theorem}
    Suppose the $V$ is a banach space and $K_1, K_2$ are two non-empty convex sets in $V$ with $K_1 \cap K_2 = \emptyset$.
    Then there exists a non-zero linear functional $f: V \to \bbR$ and a real number $\alpha$ such that
    \[f(x) \leq \alpha \leq f(y), \quad \forall x \in K_1, y \in K_2.\]
  \end{theorem}

  \begin{definition}
    Let $K$ be a convex set.
    A \emph{face} of $K$ is a convex subset $F \subseteq K$ such that for any $x, y \in K$ and any $t \in (0, 1)$, if $t x + (1 - t) y \in F$, then we have $x, y \in F$.
  \end{definition}

  \begin{definition}
    Let $K$ be a convex set and $F$ be a face of $K$.
    We say $F$ is \emph{exposed} if there exists a linear functional $f: V \to \bbR$ such that $F = \{x \in K : f(x) = \max_{y \in K} f(y)\}$.

  \end{definition}

  \begin{theorem}
    The set of exposed faces of a convex set $K$ is dense in the set of all faces of $K$.
    \Yang{}
  \end{theorem}

\subsection{Convex cones}

  \begin{definition}
    A set $C \subseteq V$ is called a \emph{convex cone} if for any $x, y \in C$ and any $\lambda, \mu \geq 0$, we have
    \[\lambda x + \mu y \in C.\]
    \Yang{}
  \end{definition}

  \begin{definition}
    Let $\calC \subset V$ be a convex cone.
    We say $\calC$ is 
     \Yang{}
  \end{definition}


\subsection{Convex bodies}

  In this subsection, we assume $V$ is an euclidean space of finite dimension $d$.

  \begin{definition}
    A convex set $K$ is called a \emph{convex body} if it is compact and has non-empty interior.
  \end{definition}

  \begin{example}[Minkowski functional]
    Let $K$ be a convex body in $V$ containing the origin in its interior.
    The \emph{Minkowski functional} of $K$ is defined as
    \[p_K(x) = \inf\{\lambda > 0 : x \in \lambda K\}.\]
    
  \end{example}

  \begin{definition}[Minkowski sum]
    Let $A_1, \cdots, A_n$ be sets in $V$ and $a_1, \cdots, a_n$ be non-negative real numbers.
    We define
    \[a_1 A_1 + \cdots + a_n A_n \coloneq \{a_1 x_1 + \cdots + a_n x_n : x_i \in A_i, i = 1, \cdots, n\}.\]
    This is called the \emph{Minkowski sum}.
  \Yang{}
  \end{definition}

  \begin{proposition}
    If $A_1, \cdots, A_n$ are convex sets (resp. convex bodies), then $a_1 A_1 + \cdots + a_n A_n$ is also a convex set (resp. convex body).
  \end{proposition}

  \begin{definition}[Haussdorff distance]
    Let $A, B$ be two non-empty compact subsets of $V$.
    The \emph{Hausdorff distance} between $A$ and $B$ is defined as
    \[d_H(A, B) = \max\{\sup_{x \in A} \inf_{y \in B} \|x - y\|, \sup_{y \in B} \inf_{x \in A} \|x - y\|\}.\]
   \Yang{}
    
  \end{definition}

  \begin{proposition}
    The Hausdorff distance $d_H$ is a metric on the set of non-empty compact subsets of $V$.

    Moreover, if $A_n \to A$ and $B_n \to B$ in the Hausdorff distance, then we have $a_1 A_n + \cdots + a_n A_n \to a_1 A + \cdots + a_n A$ in the Hausdorff distance.
   \Yang{}
  \end{proposition}

  \begin{proposition}
    The set of polytopes is dense in the set of convex bodies with respect to the Hausdorff distance.
   \Yang{}
  \end{proposition}

  We have the standard Lebesgue measure on $V$.
  Hence the volume of a convex body $K$ is well-defined, denoted by $\vol(K)$.

  \begin{theorem}
    The function $\vol(x_1 A_1 + \cdots + x_n A_n)$ is a homogeneous polynomial of degree $d$ in $x_1, \cdots, x_n$.
  \end{theorem}

  \begin{definition}
    The coefficients of the homogeneous polynomial $\vol(x_1 A_1 + \cdots + x_n A_n)$ are called the \emph{mixed volumes} of $A_1, \cdots, A_n$, denoted by $V(A_1, \cdots, A_n)$.
  \end{definition}

  \begin{proposition}
    The mixed volume $V(A_1, \cdots, A_n)$ is symmetric and multilinear in $A_1, \cdots, A_n$.

    And it is determined by the following polarization formula:
    \[V(A_1, \cdots, A_n) = \frac{1}{n!} \sum_{I \subseteq [n]} (-1)^{n - |I|} \vol\left(\sum_{i \in I} A_i\right).\]
   \Yang{}
   
  \end{proposition}

  \begin{example}[Ellipsoid]
    Let $A$ be a positive definite quadratic form on $V$.
    Then we can define an ellipsoid $E_A = \{x \in V : A(x) \leq 1\}$.
   \Yang{}
    
  \end{example}


\subsection{Positivity of mixed volumes}

  Let $V$ be an euclidean space of finite dimension $d$.
  On the set $CvBd$ of convex bodies in $V$, we have the Minkowski sum and the scalar multiplication by positive real numbers.
  One can formally define 
  \[ \Span(CvBd) = \{A - B | A, B \in CvBd\}. \]
  Let $\frakA$ be a finite dimensional subspace of $\Span(CvBd)$ such that $CvBd(\frakA) := \frakA \cap CvBd$ is a full convex cone in $\frakA$.  
  % and $\frakA = A_1, \cdots, A_n$ be convex bodies in $V$.
  Define $\upN(\frakA)$ to be the real vector space $\frakA$.
  Then we can extend the definition of mixed volumes to $\upN(\frakA)$ by multilinearity.
  Such datum $(N,N \injmap \Span(CvBd))$ is denoted by $\frakA$.

  \begin{example}[Symmetric matrices]
    Let $V$ be the space of symmetric matrices of size $d$.
    Then the set of positive definite matrices is a convex cone in $V$.
   \Yang{}
    
  \end{example}


  \begin{theorem}[Alexandrov-Fenchel inequality]
    Let $A, B, A_3, \cdots, A_n$ be convex bodies in $V$.
    Then we have
    \[V(A, A, A_3, \cdots, A_n) V(B, B, A_3, \cdots, A_n) \leq V(A, B, A_3, \cdots, A_n)^2.\]
    \Yang{} 
  \end{theorem}

  \begin{theorem}[Brunn-Minkowski inequality]
    Let $A, B$ be two non-empty compact subsets of $V$.
    Then we have
    \[\vol(A + B)^{\frac{1}{d}} \geq \vol(A)^{\frac{1}{d}} + \vol(B)^{\frac{1}{d}}.\]
    \Yang{}
  \end{theorem}

  \begin{theorem}[Minkowski inequality]
    Let $A, B$ be two convex bodies in $V$.
    Then we have
    \[V(A, A, \cdots, A)^{\frac{1}{d}} \geq V(A, A, \cdots, A, B)^{\frac{1}{d}}.\]
      \Yang{}
  \end{theorem}

